
\documentclass{article} %210 mm × 297 mm
\usepackage[utf8]{inputenc}
\usepackage[a4paper,left=1.5cm, right=1.5cm, top=2cm, bottom=2cm]{geometry}
\usepackage{graphicx}
\usepackage{systeme}          % for Solution 3
\usepackage{array}  
%\usepackage{blindtext}
\usepackage{polynom}
\usepackage{multirow}
\usepackage{mhchem}
\usepackage[normalem]{ulem}
\usepackage{url}
\usepackage{subfigure}
\usepackage{amsfonts,latexsym} % para tener disponibilidad de diversos simbolos
\usepackage{enumerate}
%\usepackage{booktabs}
\usepackage{float}
%\usepackage{threeparttable}
\usepackage{array,colortbl}
%\usepackage{ifpdf}
%\usepackage{rotating}
\usepackage{stfloats}
\usepackage{url}
\usepackage{listings}
\usepackage{fancyhdr}
\usepackage{biblatex}
\usepackage{color}
\addbibresource{sources.bib}
%\usepackage{hyperref}
%\usepackage{wrapfig}
\usepackage{array}
%\usepackage{multirow}
%\usepackage{adjustbox}
%\usepackage{nccmath}
\usepackage[dvipsnames]{xcolor}
\usepackage{tcolorbox}
\newtcolorbox{mybox}{colback=white,
colframe=black}
\newcommand{\titulo}{Abgabe 11; Diskrete Mathematik}
\newcommand{\nombre}{Liliana Sanfilippo}
\newcommand{\FCE}{\textbf{WiSe 2024/2025}}

 

\usepackage{fancyhdr}
\pagestyle{fancy}
\fancyhead{}
\fancyfoot{}
\fancyhead[LO]{\small\nombre}
\fancyhead[RE]{\small\titulo}
\fancyhead[LO]{\small\FCE}
\fancyhead[RO]{\small }
\fancyfoot[RO,LE] {\thepage}
\usepackage[onehalfspacing]{setspace}

\setcounter{page}{1} %Número de la página, la editorial lo cambiará
\begin{document}
\begin{tabular}{p{12cm}}
{{\Large\bf\titulo}}\\
\vspace{0.2cm}\\
 {\large\nombre}\\
 \vspace{0.1cm}
\end{tabular}

\section*{Präsenzaufgaben}
\begin{enumerate}
    \item Für jeden linearen Code, bestimmen Sie ob er zyklisch ist. 
    \begin{enumerate}
        \item $\{00000000, 11111111\}$,
        \begin{mybox}
        Dieser Code ist zyklisch, da alle Verschiebungen der Codewörter wieder im Code liegen. Das ist gegeben, da beide Codewörter aus nur jeweils einer Zahl bestehen. 
        \end{mybox}
        \item $\{0000, 1010, 0101, 1111\}$,
        \begin{mybox}
        Der Nullvektor und Einser-Vektor können beliebig verschoben werden und werden immer im Code bleiben. \\
        Für Verschiebungen der anderen beiden Wörter gilt: $1010 \leftrightarrow 0101$ \\
        Somit ist auch dieser Code zyklisch. 
        \end{mybox}
        \item $\{0000, 1100, 0011, 1111\}$.
         \begin{mybox}
        Verschiebt man 1100 um eine Stelle, ergibt sich 0110 und das liegt nicht im Code. Also ist der Code nicht zyklisch. 
        \end{mybox}
    \end{enumerate}
    \item Welchem Wort der Länge 4 entspricht $x^8 + x^5 + x^3 + x + 1 \in \mathbb{Z}/2\mathbb{Z}[x]$?
    \begin{mybox}
    Da gilt $x^{n} = x$, gilt $x^{4+1} = x^1$ und $x^{4+4} = x^4 = x$
    \[
    x^8 + x^5 + x^3 + x + 1 \equiv x + x^2 + x^3 + x + 1 = 1 + x^2 + x^3
    \]
    Das entsprechende Wort ist 1011.
    \end{mybox}
\end{enumerate}
\begin{mybox}
\textbf{Verständnisfrage}: \\
Muss man immer von dem kleinsten Grad zum größten lesen? 
Also $1 + x^2 + x^3$ ist 1011 und man kann es nicht als 1101 lesen?
\end{mybox}

\section*{Aufgaben zur Abgabe}
\subsection*{Aufgabe 1 (4 Punkte)}  
Finden Sie alle Polynome von Grad höchstens 5 im Ideal  
$\langle x^3 + x \rangle \subseteq \mathbb{Z}/2\mathbb{Z}[x]$.
\begin{mybox}
Damit Polynome $(x^3 + x) \cdot q(x)$ höchstens den Grad 5 haben, darf $q(x)$ höchstens den Grad 2 haben, da $x^3 + x$ den Grad 3 hat. \\
Also $q(x)=a+bx+cx^2$ mit $a, b, c \in \mathbb{Z}/2\mathbb{Z}$
\begin{table}[H]
    \centering
    \[
    \begin{tabular}{c|c}
        $q(x)$ & $q(x)g(x)$  (Vielfaches von $x^3 + x)$   \\ \hline
        $0$ & $0$\\
        $1$ & $x+x^3$  \\
        $x$ & $x^2+x^4$  \\
        $1+x$ & $x^4 + x^3 + x^2 + x$ \\
        $x^2$ & $x^5 + x^3$  \\
        $1+x^2$ & $x^5  + x $ \\
        $x+x^2$ & $x^5 + x^4 + x^3 + x^2$ \\
        $1+x+x^2$ & $x^5 + x^4  + x^2 + x$  
    \end{tabular}
    \]
\end{table}
\end{mybox}
\subsection*{Aufgabe 2 (4 Punkte)}  
Zeigen Sie, dass  
\[ 
(x + 1)^7 = x^7 + 1 \text{ in } \mathbb{Z}/7\mathbb{Z}[x].
\]  
\begin{mybox}
%Wir müssen alles modulo 7 rechnen. 
%\begin{align*}
%    (x + 1)^7 \mod 7 & = (x^7 + 7 x^6 + 21 x^5 + 35 x^4 + 35 x^3 + 21 x^2 + 7 x + 1) \mod 7 \\ 
%    & = x^7 + x^6 + 3 x^5 + 5 x^4 + 5 x^3 + 3 x^2 + x + 1 
%\end{align*}
Der binomische Lehrsatz ergibt in diesem Fall: 
\begin{align*}
    (x+1)^7 & = \sum_{k=0}^7 \binom{7}{k} x^{7-k} \cdot 1^k  \mod 7 \\
    & = \sum_{k=0}^7 \binom{7}{k} x^{7-k} \mod 7 \\
    & = \binom{7}{0} x^{7} + \binom{7}{1} x^{6} + \binom{7}{2} x^{5} + \binom{7}{3} x^{4} + \binom{7}{4} x^{3} + \binom{7}{4} x^{3} + \binom{7}{5} x^{2} + \binom{7}{6} x^{1} + \binom{7}{7} x^{0} \mod 7 \\
     & = 1 x^{7} + \binom{7}{1} x^{6} + \binom{7}{2} x^{5} + \binom{7}{3} x^{4} + \binom{7}{4} x^{3} + \binom{7}{4} x^{3} + \binom{7}{5} x^{2} + \binom{7}{6} x^{1} + 1 x^{0} \mod 7 \\ 
     & = 1 x^{7}  + 1 x^{0} \\ 
     & = x^{7}  +  x^{0} \\ 
     & = x^{7}  +  1 
\end{align*}

\end{mybox}
Für welche positiven ganzen Zahlen $m$ gilt $(x + 1)^m = x^m + 1$ in $\mathbb{Z}/m\mathbb{Z}[x]$?

\subsection*{Aufgabe 3 (4 Punkte)}  
Bestimmen Sie den zyklischen Code $C$ der Länge 3 mit dem Erzeugendenpolynom  
$g(x) = x^2 + x + 1 \in \mathbb{Z}/2\mathbb{Z}[x]$.

\begin{mybox}
Weil der Code die Länge 3 hat, muss ich $p(x)g(x) \mod (x^3 + 1)$ berechnen. 
\begin{table}[H]
    \centering
   \[
    \begin{tabular}{c|c|c|c}
        $p(x)$ & $p(x)g(x)$  &   $p(x)g(x) \mod (x^3 + 1)$ & \text{Codewort}  \\ \hline
        $0$ & $0$& $0$ & 000 \\
        $1$ & $1+x+x^2$ & $1+x+x^2$ & 111 \\
        $x$ & $x+x^2+x^3$ & $x+x^2$ & 011 \\
        $1+x$ & $x^3  + 1$ & $0$ & 000 \\
        $x^2$ & $x + x^3 + x^2$ & $x + x^2$ & 011 \\
        $1+x^2$ & $x^3  + 1 $ & $0$ & 000 \\
        $x+x^2$ & $0$ & $0$ & 000\\
        $1+x+x^2$ & $  x^2 +  x + 1$ &$  x^2 +  x + 1$ & 111 \\
    \end{tabular}
   \]
\end{table}
Der Code ist also: \{000, 011, 111\}
\end{mybox}

\begin{mybox}
\textbf{Verständnisfrage}: \\
Bei $p(x) = 1+x^2$ hatte ich folgenden Gedankenweg: 
\begin{align*}
p(x)g(x)& =   (1+x^2) \cdot (x^2 + x + 1) \\
        & \neq x^4 + x^3 + 2 x^2 + x + 1 \mod 2 & \text{nicht, da es ja ein Restklassenring mod 2 ist} \\
        & = x^1 + x^3 + 2 x^2 + x + 1 \mod 2 & \\
         & =   x^3 + 2 x^2 + 2 x + 1 \mod 2 & \text{zusammengerechnet} \\
        & =  x^3  + 1  
\end{align*}
Ist das so korrekt nachvollzogen? 
\end{mybox}

\subsection*{Aufgabe 4 (4 Punkte)}  
Bestimmen Sie den zyklischen Code $C$ der Länge 3 mit dem Kontrollpolynom  
$h(x) = x^2 + x + 1 \in \mathbb{Z}/2\mathbb{Z}[x]$.
\begin{mybox}
Wenn $x^2 + x + 1$ das Kontrollpolynom ist und wir den Code $C$ bestimmen wollen, müssen wir aus $h(x)$ das Erzeugendenpolynom $g(x)$ ableiten und daraus den Code erzeugen. \\
$x^n + 1 = g(x)h(x)$, also $g(x) = \frac{x^n + 1}{h(x)}$ und hier $g(x) = \frac{x^3 + 1}{x^2 + x + 1}$
\[
\begin{array}{cccccccccccc}
    & x^3   &       &       & +1 &  \\
-(  & x^3   & + x^2 & + x   &    &) & | & h \text{ multipliziert mit } x\\\hline
    &       & + x^2 & + x   & +1 & \\
-(  &       & + x^2 & + x   & +1 &) & | & h \text{ multipliziert mit } 1 \\\hline
& & & & 0
\end{array}
\]
Also gibt es keinen Rest und $q(x) = x+1$ \\
Aus dem Skript wissen wir, dass mit $n = 3$ und $g(x) = x+1$ ein Code $ C = \{000, 110, 011, 101\}$ ist. 
\end{mybox}
\end{document}
